\documentclass[12pt]{article}
\PassOptionsToPackage{authoryear, compress}{natbib}

\usepackage[utf8]{inputenc}
\usepackage[T1]{fontenc}
\usepackage{hyperref}
\usepackage{url}
\usepackage{booktabs}
\usepackage{amsfonts}
\usepackage{nicefrac}
\usepackage{microtype}
\usepackage{xcolor}
\usepackage[authoryear]{natbib}
\usepackage{graphicx}
\usepackage{subcaption}
\usepackage{multirow}
\usepackage{placeins}
\usepackage{amsmath}
\usepackage[margin=1.2in]{geometry}

\title{Uncertainty-Aware Reconstruction of Airfoil Surface Pressure\\
       from Sparse Tap Measurements via Conditional Flow Matching}

\author{Atharva Mahajan \\
        University of Michigan \\
        \texttt{atharvm@umich.edu}}

\begin{document}

\maketitle

\begin{abstract}
Recovering full surface pressure distributions from sparse sensor measurements
is a fundamental challenge in experimental aerodynamics. Dense tap arrays and
Pressure-Sensitive Paint are expensive and often geometrically infeasible,
motivating data-driven reconstruction from limited observations. We present a
generative framework based on \textbf{Conditional Flow Matching (CFM)} that
reconstructs the complete 2D airfoil pressure field from $k \in \{5, 10, 20\}$
pressure tap measurements, while providing calibrated uncertainty estimates over
the inferred field. The model combines a \textbf{DeepSets encoder with learned
positional attention} — which aggregates sparse, arbitrarily-placed tap
observations into a compact conditioning vector — with a \textbf{1D convolutional
UNet vector field} conditioned via Feature-wise Linear Modulation (FiLM) at
every spatial scale. Trained on the AirfRANS dataset of 800 RANS simulations,
the framework achieves a mean absolute error of 183\,Pa (14\% of the pressure
standard deviation) at $k=10$ uniformly-spaced taps, and recovers lift and drag
coefficients to within $\Delta C_L = 0.094$ and $\Delta C_D = 0.015$
respectively. We further evaluate a randomized tap placement variant that
generalises to arbitrary sensor configurations, achieving 237\,Pa MAE at
$k=20$ — providing a quantitative bound on the cost of not controlling sensor
placement. The posterior ensemble is well-calibrated at $k=5$ (97\% coverage
at $\pm2\sigma$) and remains conservative at higher tap counts, making the
framework suitable for risk-aware aerodynamic analysis.
\end{abstract}

%------------------------------------------------------------------------------
\section{Introduction}
\label{sec:intro}
%------------------------------------------------------------------------------

The surface pressure distribution is the primary aerodynamic quantity governing
lift, drag, and pitching moment. In flight testing and wind tunnel experiments,
it is measured by arrays of pressure taps — small orifices connected to
transducers — whose placement is constrained by structural access, cost, and
the geometry of the model. As a result, experimental campaigns routinely operate
with $\mathcal{O}(10$--$100)$ taps on surfaces with $\mathcal{O}(10^3)$
meaningful spatial degrees of freedom, requiring interpolation or model-based
reconstruction to recover the full field.

Classical reconstruction approaches include polynomial fitting, radial basis
function interpolation, and physics-informed splines. These methods assume
smoothness but do not quantify uncertainty, and degrade rapidly when taps are
sparse relative to the complexity of the pressure distribution — particularly
near the leading-edge suction peak and in regions of flow separation.

Generative deep learning offers a natural alternative: rather than producing a
single reconstruction, a generative model learns the conditional distribution
$p(\mathbf{p} \mid \mathbf{y}_k)$, where $\mathbf{p} \in \mathbb{R}^N$ is the
full pressure field and $\mathbf{y}_k$ are the $k$ tap observations. Drawing
samples from this posterior provides both a mean reconstruction and uncertainty
estimates that reflect genuine ambiguity in the under-determined problem.

\textbf{Flow Matching} \citep{lipman2022flow} is a recently proposed generative
paradigm that learns a deterministic ODE-based transport from a Gaussian
reference to the target distribution, trained by regressing onto the
ground-truth velocity field along straight-line probability paths. Compared to
score-based diffusion models, it offers faster training convergence, simpler
implementation, and more stable inference trajectories. \textbf{Conditional Flow
Matching} (CFM) \citep{albergo2022building} extends this to conditional generation
by incorporating an observation vector into the learned velocity field.

This work applies CFM to 2D airfoil pressure reconstruction. Our contributions
are:
\begin{enumerate}
    \item A complete CFM pipeline for sparse aerodynamic field reconstruction,
          integrating a DeepSets encoder with learned positional attention and a
          1D conv UNet vector field with FiLM conditioning.
    \item A controlled ablation of uniform vs.\ randomized tap placement,
          quantifying the reconstruction cost of unconstrained sensor locations.
    \item A detailed analysis of model failure modes by angle of attack and
          flow regime, identifying the operating envelope where the framework
          is reliable.
\end{enumerate}

%------------------------------------------------------------------------------
\section{Related Work}
\label{sec:related}
%------------------------------------------------------------------------------

\paragraph{Sparse field reconstruction.}
Sparse sensor placement and reconstruction has been studied in the context of
gappy POD \citep{everson1995karhunen}, optimal sensor placement
\citep{manohar2018data}, and compressed sensing \citep{bai2020low}. These methods
require a basis computed from the training ensemble and cannot extrapolate to
conditions outside it. Neural operator approaches \citep{li2021fourier} learn
mappings between function spaces but typically require dense input observations.

\paragraph{Generative models for physical fields.}
Score-based diffusion has been applied to turbulence super-resolution
\citep{shu2023physics} and weather reconstruction \citep{price2023gencast}.
Flow Matching has been adopted for molecular generation and video synthesis;
its application to physical fields is nascent. Closest to our work is
Diff-SPORT, which applies diffusion to sparse sports tracking, demonstrating
that conditional score models generalise well to irregular observation sets.

\paragraph{Aerodynamic surrogate modelling.}
Deep learning surrogates for airfoil aerodynamics have been explored extensively
\citep{thuerey2020deep, bhatnagar2019prediction}, primarily for predicting scalar
coefficients or full-field volume quantities. Surface pressure reconstruction
from sparse taps has received less attention; \citet{lee2020deep} studied CNN
interpolation on structured grids, but without uncertainty quantification or
support for irregular tap placements.

%------------------------------------------------------------------------------
\section{Problem Formulation}
\label{sec:problem}
%------------------------------------------------------------------------------

Let $\mathbf{p} \in \mathbb{R}^N$ denote the surface pressure field sampled at
$N=256$ uniformly-spaced arc-length positions $s_i \in [0,1]$ around the airfoil,
where $s=0$ is the leading edge and $s=1$ closes back at the trailing edge.

A sparse observation set $\mathbf{y}_k = \{(s_j, p_j)\}_{j=1}^k$ consists of
$k$ pressure measurements at tap positions $\{s_j\}$. The inlet conditions
$\mathbf{u} = [v_x, v_y, \alpha]$ — freestream velocity components and angle of
attack — are assumed known. The reconstruction task is to infer the posterior
$p(\mathbf{p} \mid \mathbf{y}_k, \mathbf{u})$.

In the \textbf{uniform} tap setting, tap positions are fixed at
$s_j = (j-1)/k$ for $j = 1, \ldots, k$, providing even surface coverage.
In the \textbf{randomized} setting, tap indices are drawn uniformly at random
from $\{0, \ldots, N-1\}$ without replacement, sorted by arc position before
encoding, reflecting realistic experimental constraints.

%------------------------------------------------------------------------------
\section{Methodology}
\label{sec:method}
%------------------------------------------------------------------------------

\subsection{Conditional Flow Matching}

We adopt the straight-line OT-CFM formulation \citep{lipman2022flow}. Given a
target pressure field $\mathbf{x}_1 \sim p_{\text{data}}$ and Gaussian noise
$\mathbf{x}_0 \sim \mathcal{N}(0, I)$, we define the interpolant:
\begin{equation}
    \mathbf{x}_t = (1 - (1 - \sigma_\mathrm{min})\,t)\,\mathbf{x}_0 + t\,\mathbf{x}_1,
    \quad t \in [0, 1],
    \label{eq:interpolant}
\end{equation}
with $\sigma_\mathrm{min} = 10^{-4}$. The target velocity along this path is:
\begin{equation}
    \mathbf{u}_t = \mathbf{x}_1 - (1 - \sigma_\mathrm{min})\,\mathbf{x}_0.
\end{equation}
The velocity network $v_\theta(\mathbf{x}_t, t, \mathbf{c})$ minimises:
\begin{equation}
    \mathcal{L}_{\mathrm{CFM}}(\theta) =
    \mathbb{E}_{t,\, \mathbf{x}_1,\, \mathbf{x}_0}
    \left[\left\| v_\theta(\mathbf{x}_t, t, \mathbf{c}) - \mathbf{u}_t \right\|^2\right].
    \label{eq:cfm_loss}
\end{equation}
At inference, $n_\mathrm{samples}$ noise vectors are integrated forward via
Euler's method (20 steps), producing a posterior ensemble over pressure fields.

\subsection{Sparse Tap Encoder}

The condition vector $\mathbf{c} \in \mathbb{R}^{128}$ is computed by a
\textbf{DeepSets} encoder with learned positional attention:
\begin{equation}
    \mathbf{c} = f_\mathrm{fuse}\!\left(\textstyle\sum_{j=1}^k w_j\, g(\mathbf{y}_j),\;\mathbf{u},\; k/k_\mathrm{max}\right),
    \label{eq:encoder}
\end{equation}
where $g$ is a shared per-tap MLP, $w_j = \mathrm{softmax}(a(s_j))_j$ are
attention weights from a position MLP $a$, and $f_\mathrm{fuse}$ is a fusion
MLP. Table~\ref{tab:encoder} details each component.

\begin{table}[ht]
\centering
\caption{Sparse tap encoder architecture.}
\label{tab:encoder}
\begin{tabular}{lllc}
\toprule
Component & Input & Output & Layers \\
\midrule
Per-tap MLP $g$ & $(s_j, p_j) \in \mathbb{R}^2$ & $\mathbb{R}^{128}$ & 3$\times$Linear(128) + SiLU \\
Attention MLP $a$ & $s_j \in \mathbb{R}$ & scalar logit & Linear(32) + SiLU + Linear(1) \\
Softmax pooling & $k$ embeddings + logits & $\mathbb{R}^{128}$ & $\sum_j \mathrm{softmax}(a_j)\, g_j$ \\
Fusion MLP $f_\mathrm{fuse}$ & $[\mathbf{h};\,\mathbf{u};\,k/k_\mathrm{max}] \in \mathbb{R}^{132}$ & $\mathbf{c} \in \mathbb{R}^{128}$ & 2$\times$Linear(128) + SiLU \\
\bottomrule
\end{tabular}
\end{table}

The attention weight $a(s_j)$ allows the model to learn that taps at physically
significant positions — the suction peak ($s \approx 0.05$), trailing edge
($s \approx 0$) — are more informative than taps on the flat pressure side.
Including $k/k_\mathrm{max}$ in the fusion enables the vector field to produce
sharper posteriors when more taps are available. Under uniform placement,
attention weights converge to position-independent constants (mean pooling is
equivalent); under random placement the attention is genuinely informative, as
the model must discover which positions are most constraining from the pressure
values alone.

\textbf{Why not a convolutional encoder?} With $k \leq 20$ taps on a 256-point
grid, the observation signal is 92--98\% zeros when scattered onto the grid.
A 1D conv encoder learns to detect isolated impulses rather than spatial context,
empirically underperforming the DeepSets baseline. Conv inductive biases benefit
dense signals — motivating their use in the vector field but not the encoder.

\subsection{1D Convolutional UNet Vector Field}

The velocity network $v_\theta$ is a 1D UNet (Table~\ref{tab:unet}) operating
on the 256-point pressure signal. Each residual block uses
\textbf{FiLM conditioning}:
\begin{equation}
    \mathbf{h}' = \mathbf{h} + \sigma\!\left(\mathrm{GN}(\mathrm{Conv}(\mathbf{h})) \cdot (1 + \boldsymbol{\gamma}) + \boldsymbol{\beta}\right),
    \label{eq:film}
\end{equation}
where $\boldsymbol{\gamma}, \boldsymbol{\beta}$ are linear projections of
$[\mathbf{c};\,\mathbf{t}_\mathrm{emb}]$ and $\mathbf{t}_\mathrm{emb} \in
\mathbb{R}^{32}$ is a Fourier time embedding.

\begin{table}[ht]
\centering
\caption{1D conv UNet vector field architecture. Residual block convolutions use
circular padding; strided down/upsample convolutions use zero padding. GN =
GroupNorm(8). Each level contains $n_\mathrm{blocks}=2$ FiLM residual blocks.
Skip connections link encoder to decoder at each level.}
\label{tab:unet}
\begin{tabular}{lcccc}
\toprule
Level & Spatial size & Channels & Operation & Kernel \\
\midrule
Input proj.   & 256 & 1 $\to$ 64   & Conv1d(1)    & 1 \\
Encoder 0     & 256 & 64           & 2$\times$ResBlock & 5 \\
Downsample 0  & 256$\to$128 & 64$\to$128  & Conv1d(stride 2) & 4 \\
Encoder 1     & 128 & 128          & 2$\times$ResBlock & 5 \\
Downsample 1  & 128$\to$64  & 128$\to$256 & Conv1d(stride 2) & 4 \\
Bottleneck    & 64  & 256          & 2$\times$ResBlock & 5 \\
Upsample 1    & 64$\to$128  & 256$\to$128 & ConvTranspose1d  & 4 \\
Decoder 1     & 128 & 128 (+skip)  & 2$\times$ResBlock & 5 \\
Upsample 0    & 128$\to$256 & 128$\to$64  & ConvTranspose1d  & 4 \\
Decoder 0     & 256 & 64 (+skip)   & 2$\times$ResBlock & 5 \\
Output proj.  & 256 & 64$\to$1     & Conv1d(1)    & 1 \\
\bottomrule
\end{tabular}
\end{table}

\textbf{Why FiLM at every layer?} In preliminary experiments, concatenating
$\mathbf{c}$ only at the input caused posterior collapse: the network ignored
the condition and predicted the marginal pressure distribution (val loss
$\approx 0.36$, MAE $\approx 250$\,Pa). FiLM injects the condition at every
block, making it structurally impossible to deactivate the conditioning signal.

\textbf{Circular padding in residual blocks.} The residual convolutions
use circular padding so that the receptive field wraps correctly at the
arc-length endpoints — $s=0$ and $s=1$ are the same physical point (trailing
edge). Strided down/upsample convolutions use standard zero padding, as they
operate on compressed representations where boundary effects are less significant.

\subsection{Integrated Aerodynamic Coefficients}

Lift and drag coefficients are computed from each ensemble member via:
\begin{equation}
    C_L = \frac{1}{q_\infty}\int_0^L p(s)\,\mathbf{n}(s)\cdot\hat{\mathbf{l}}\,ds, \qquad
    C_D = \frac{1}{q_\infty}\int_0^L p(s)\,\mathbf{n}(s)\cdot\hat{\mathbf{d}}\,ds,
\end{equation}
evaluated by the trapezoidal rule on the 256-point arc-length grid, where
$q_\infty = \tfrac{1}{2}\|v\|^2$ is the dynamic pressure.

\subsection{Training Details}

Pressure fields are standardised using training statistics
($\mu_p = -488$\,Pa, $\sigma_p = 1330$\,Pa). The model is trained for 400
epochs with AdamW (lr $= 3\times10^{-4}$, weight decay $= 10^{-4}$), cosine
annealing, and gradient clipping (norm 1.0). Total parameters: 1.92\,M.
Training takes $\approx$10 minutes on a single NVIDIA GPU.

%------------------------------------------------------------------------------
\section{Experimental Setup}
\label{sec:experiments}
%------------------------------------------------------------------------------

\subsection{Dataset}

We use the \textbf{AirfRANS} dataset \citep{bonnet2022airfrans}, comprising
1,000 RANS simulations of 2D airfoils at angles of attack $-5^\circ$ to
$+15^\circ$ and freestream speeds 31--94\,m/s. The surface boundary is extracted
by tracing the VTP connectivity graph from the leading edge and resampled to a
uniform 256-point arc-length grid. A stratified 80/20 split by $C_D$ decile
ensures matched flow-regime distributions in train and test sets (an unconstrained
random split produced $1.7\times$ more stalled cases in the test set).

\subsection{Evaluation Metrics}

For each test case, we draw $n_\mathrm{samples}=100$ posterior samples and report:
\textbf{MAE($p$)} in Pa; \textbf{Relative MAE} as a fraction of the per-case
pressure range; \textbf{MAE($C_L$)}, \textbf{MAE($C_D$)} from the ensemble mean;
and \textbf{coverage} — the fraction of cases where the ground truth falls within
$\pm2\sigma$ of the posterior (target: 95\%).

%------------------------------------------------------------------------------
\section{Results}
\label{sec:results}
%------------------------------------------------------------------------------

\subsection{Qualitative Reconstruction}

Figures~\ref{fig:uniform_recon} and~\ref{fig:random_recon} show reconstruction
results for a representative test case under uniform and randomized tap placement
respectively. Each figure shows, for $k \in \{5, 10, 20\}$: (top) the airfoil
geometry with numbered tap locations coloured by their measured pressure, and
(bottom) the posterior mean $\pm1\sigma$/$\pm2\sigma$ bands vs.\ ground truth.

Under uniform placement (Figure~\ref{fig:uniform_recon}), the model reliably
captures the leading-edge suction peak, mid-chord pressure recovery, and
trailing-edge stagnation. Uncertainty bands narrow visibly from $k=5$ to $k=10$.
The slight spike at the trailing edge ($s \approx 0.5$) is a real physical
discontinuity — the pressure jump between suction and pressure surfaces at the
trailing edge — and is not a reconstruction artefact.

Under randomized placement (Figure~\ref{fig:random_recon}), reconstruction
quality is more variable. When random taps happen to land near the suction peak
(observable from the numbered locations on the airfoil), the reconstruction is
comparable to uniform. When taps cluster on the flat pressure side, the suction
peak is poorly constrained and uncertainty bands remain wide. This case-by-case
variability is a fundamental property of random placement, and is correctly
reflected in the wider uncertainty bands.

\begin{figure}[ht]
    \centering
    \includegraphics[width=\textwidth]{figures_uniform/fig1.png}
    \caption{Uniform tap placement. \textit{Top row}: airfoil geometry with
             numbered tap locations (coloured by measured pressure). \textit{Bottom
             row}: posterior pressure reconstruction. Solid: ensemble mean. Bands:
             $\pm1\sigma$ (dark), $\pm2\sigma$ (light). Dashed: ground truth.
             Numbers on the pressure plot correspond to tap labels on the airfoil.}
    \label{fig:uniform_recon}
\end{figure}

\begin{figure}[ht]
    \centering
    \includegraphics[width=\textwidth]{figures_random/fig1.png}
    \caption{Randomized tap placement for the same test case. Tap locations vary
             stochastically; when taps cluster away from the suction peak
             (e.g.\ $k=5$), uncertainty bands are wider and the peak is less
             accurately resolved.}
    \label{fig:random_recon}
\end{figure}

\subsection{Overall Performance}

Table~\ref{tab:main_results} reports aggregate performance across 100 test cases.

\begin{table}[ht]
\centering
\caption{Reconstruction performance (100 test cases, 100 samples each).
Coverage target: 95\% at $\pm 2\sigma$.}
\label{tab:main_results}
\begin{tabular}{llccccc}
\toprule
Strategy & $k$ & MAE($p$) [Pa] & MAE($C_L$) & MAE($C_D$) & $C_L$ cov. & $C_D$ cov. \\
\midrule
\multirow{3}{*}{Uniform}
  & 5  & 211 & 0.115 & 0.017 & 97\% & 97\% \\
  & 10 & 183 & 0.094 & 0.015 & 100\% & 98\% \\
  & 20 & 189 & 0.094 & 0.015 & 100\% & 98\% \\
\midrule
\multirow{3}{*}{Random}
  & 5  & 275 & 0.149 & 0.022 & 96\% & 95\% \\
  & 10 & 255 & 0.134 & 0.020 & 96\% & 96\% \\
  & 20 & 237 & 0.122 & 0.019 & 100\% & 97\% \\
\bottomrule
\end{tabular}
\end{table}

Uniform placement outperforms random by 25--35\% in MAE($p$). The two strategies
show contrasting $k$-scaling: uniform saturates (and slightly regresses) between
$k=10$ and $k=20$ due to discrete tap-stride alignment effects, while random
placement improves monotonically as larger $k$ stochastically guarantees better
coverage of high-gradient regions. $C_L$ and $C_D$ errors scale sublinearly with
pressure MAE, as coefficient errors depend on surface integrals that average out
local reconstruction noise.

\subsection{Aerodynamic Coefficient Posteriors}

Figures~\ref{fig:uniform_clcd} and~\ref{fig:random_clcd} show the posterior
distributions of $C_L$ and $C_D$ as violin plots for a representative test case.
Each violin represents the ensemble of 200 samples drawn from the CFM posterior
at a given $k$, with the dashed line marking the ground truth.

Under uniform placement (Figure~\ref{fig:uniform_clcd}), the $C_L$ posterior
narrows and converges toward the ground truth as $k$ increases, reflecting the
additional information carried by more taps. The $C_D$ posterior similarly
tightens but retains meaningful spread due to the inherent sensitivity of drag
to local pressure features near the trailing edge.

Under randomized placement (Figure~\ref{fig:random_clcd}), the posteriors are
wider at all $k$, consistent with the higher pressure MAE in Table~\ref{tab:main_results}.
The ground truth is contained within the posterior in all cases, confirming
that the model does not produce overconfident estimates even under random placement
— consistent with the over-coverage observed in Table~\ref{tab:main_results}.

\begin{figure}[ht]
    \centering
    \includegraphics[width=0.85\textwidth]{figures_uniform/fig3.png}
    \caption{$C_L$ and $C_D$ posterior distributions for $k \in \{5, 10, 20\}$
             uniform taps. Violin width represents posterior density. Dashed line:
             ground truth. The posterior narrows with increasing $k$, correctly
             reflecting reduced uncertainty from more observations.}
    \label{fig:uniform_clcd}
\end{figure}

\begin{figure}[ht]
    \centering
    \includegraphics[width=0.85\textwidth]{figures_random/fig3.png}
    \caption{$C_L$ and $C_D$ posteriors for randomized tap placement. Wider
             distributions reflect the additional uncertainty from not controlling
             sensor locations. The ground truth remains within the posterior
             at all $k$, confirming conservative but not overconfident uncertainty
             estimates.}
    \label{fig:random_clcd}
\end{figure}

\subsection{Failure Mode Analysis by Angle of Attack}

Figure~\ref{fig:uniform_aoa} and~\ref{fig:random_aoa} show per-case MAE vs.\
angle of attack across the full test set, with bin means overlaid. Both absolute
MAE (Pa) and relative MAE (as a fraction of the per-case pressure range) are
shown; the relative metric removes the confound of high-speed cases having larger
absolute pressure magnitudes. Table~\ref{tab:aoa_breakdown} summarises the AoA
breakdown for the uniform model.

\begin{table}[ht]
\centering
\caption{MAE($p$) by AoA bin, uniform placement, $k=10$.}
\label{tab:aoa_breakdown}
\begin{tabular}{lccc}
\toprule
AoA range & $n$ & Mean MAE [Pa] & 90th pct [Pa] \\
\midrule
$-5^\circ$ to $0^\circ$    & 54 & 219 & 553 \\
$0^\circ$  to $+5^\circ$   & 60 & 137 & 261 \\
$+5^\circ$ to $+10^\circ$  & 40 & 102 & 164 \\
$+10^\circ$ to $+15^\circ$ & 46 & 109 & 210 \\
\bottomrule
\end{tabular}
\end{table}

Negative AoA cases are systematically harder: all 10 worst-performing cases occur
at $-2^\circ$ to $-5^\circ$. At negative incidence the suction peak migrates to
the lower surface, appearing at an unfamiliar arc-length position relative to the
positive-AoA training distribution. Both strategies show this failure mode, though
it is more pronounced under random placement (Figure~\ref{fig:random_aoa}) where
the suction peak is also less likely to be directly sampled. Notably, the relative
MAE plot shows that negative-AoA cases are harder even after normalising by
pressure range, confirming the failure is structural rather than a consequence
of lower dynamic pressure.

\begin{figure}[ht]
    \centering
    \includegraphics[width=\textwidth]{figures_uniform/fig4.png}
    \caption{Per-case MAE vs.\ AoA, uniform taps, $k=10$. \textit{Left}: absolute
             MAE (Pa). \textit{Right}: relative MAE as a percentage of the per-case
             pressure range, removing the effect of varying dynamic pressure.
             Scatter: individual test cases. Red line: bin means. Degradation below
             $0^\circ$ is visible in both metrics, confirming a structural failure
             mode at negative incidence.}
    \label{fig:uniform_aoa}
\end{figure}

\begin{figure}[ht]
    \centering
    \includegraphics[width=\textwidth]{figures_random/fig4.png}
    \caption{Per-case MAE vs.\ AoA for randomized taps, $k=10$. The negative-AoA
             failure mode is more pronounced: random placement reduces the
             probability of sampling the suction peak at any incidence, and
             at negative AoA the peak position is also unfamiliar to the model.}
    \label{fig:random_aoa}
\end{figure}

Table~\ref{tab:cd_breakdown} shows MAE by $C_D$ regime. Counterintuitively,
stalled cases (CD $> 0.05$) are easier to reconstruct than mildly separated ones:
stalled airfoils have broad, featureless low-pressure plateaux that sparse taps
can capture, whereas mildly separated cases have sharp local features (separation
bubbles, reattachment) sensitive to exact tap placement.

\begin{table}[ht]
\centering
\caption{MAE($p$) by $C_D$ regime, uniform placement, $k=10$.}
\label{tab:cd_breakdown}
\begin{tabular}{lccc}
\toprule
$C_D$ range & Flow regime & $n$ & Mean MAE [Pa] \\
\midrule
$[0.00, 0.01)$ & Fully attached  & 119 & 150 \\
$[0.01, 0.05)$ & Mild separation &  30 & 183 \\
$[0.05, 0.10)$ & Separated       &  15 &  99 \\
$[0.10, 0.25]$ & Stalled         &  31 & 119 \\
\bottomrule
\end{tabular}
\end{table}

%------------------------------------------------------------------------------
\section{Discussion}
\label{sec:discussion}
%------------------------------------------------------------------------------

\paragraph{Uniform vs.\ randomized placement.}
The 25--35\% MAE advantage of uniform placement quantifies a practically important
tradeoff. For new test articles where tap locations can be specified, even spacing
is significantly more informative per tap and provides better-behaved $k$-scaling.
For legacy models with fixed instrumentation, the randomized model provides a
viable reconstruction path, with the attention mechanism learning to exploit
whichever taps happen to land near informative regions.

\paragraph{Contrasting $k$-scaling.}
Uniform placement saturates between $k=10$ and $k=20$ (183 $\to$ 189\,Pa) due
to discrete stride effects: the $k=20$ tap set is not a superset of $k=10$, so
some cases lose coverage of the suction peak. Random placement improves
monotonically (275 $\to$ 255 $\to$ 237\,Pa), as larger $k$ stochastically
guarantees better high-gradient coverage in expectation. This suggests random
placement offers a more predictable sensor-count scaling law.

\paragraph{Limitations.}
The model is trained exclusively on SST RANS simulations from AirfRANS, limiting
validity to that distribution. The negative AoA regime ($-5^\circ$ to $0^\circ$)
is a known weak point. The over-dispersed posterior at $k > 5$ suggests the CFM
objective does not fully exploit additional tap information; tighter calibration
could be achieved via classifier-free guidance, at the cost of training complexity.

%------------------------------------------------------------------------------
\section{Conclusion}
\label{sec:conclusion}
%------------------------------------------------------------------------------

We presented a Conditional Flow Matching framework for uncertainty-aware
reconstruction of airfoil surface pressure distributions from sparse tap
measurements. The model combines a DeepSets encoder with learned positional
attention — enabling generalisation to arbitrary sensor placements — with a
1D convolutional UNet vector field conditioned via FiLM at every spatial scale.
On the AirfRANS benchmark, the framework achieves 183\,Pa MAE at $k=10$ uniform
taps with well-calibrated uncertainty. The controlled ablation against randomized
placement quantifies the 25--35\% reconstruction cost of not controlling sensor
location, and reveals contrasting $k$-scaling behaviour that has practical
implications for experimental sensor design.

Future work will explore adaptive tap placement via active learning, extension to
transonic regimes using the UniFoil dataset, and application to 3D wing pressure
reconstruction from surface tap arrays.

\newpage
\bibliographystyle{abbrv}
\bibliography{ref_report}

\end{document}
